\documentclass[10pt,a4paper]{article}
\usepackage[utf8]{inputenc}
\usepackage[T1]{fontenc}
\usepackage[spanish]{babel}
\usepackage{amsmath,amssymb}
\usepackage[margin=1.2cm]{geometry}
\usepackage{multicol}
\usepackage{array}

\setlength{\parindent}{0pt}
\setlength{\columnsep}{0.9cm}
\pagestyle{empty}

% Compactar el espaciado alrededor de las ecuaciones
\setlength{\abovedisplayskip}{3pt}
\setlength{\belowdisplayskip}{3pt}
\setlength{\abovedisplayshortskip}{2pt}
\setlength{\belowdisplayshortskip}{2pt}

\newcommand{\seccion}[1]{%
  \vspace{2mm}\hrule\vspace{1mm}%
  \textbf{\footnotesize\MakeUppercase{#1}}\vspace{1mm}\hrule\vspace{2mm}}

\begin{document}
\small

\begin{center}
  \textbf{\large FORMULARIO — VIBRACIONES MECÁNICAS}
\end{center}
\vspace{-2mm}

\begin{multicols}{2}

\seccion{Movimiento armónico simple}

Ecuación diferencial del sistema libre no amortiguado:
\begin{equation*}
  m\ddot{x} + kx = 0
  \qquad\Longleftrightarrow\qquad
  \ddot{x} + \omega_n^{2}x = 0
\end{equation*}

Ecuación general del movimiento:
\begin{equation*}
  x(t) = A\cos(\omega_n t) + B\sin(\omega_n t)
\end{equation*}

Primera derivada (velocidad):
\begin{equation*}
  \dot{x}(t) = -A\,\omega_n\sin(\omega_n t) + B\,\omega_n\cos(\omega_n t)
\end{equation*}

Segunda derivada (aceleración):
\begin{equation*}
  \ddot{x}(t) = -A\,\omega_n^{2}\cos(\omega_n t) - B\,\omega_n^{2}\sin(\omega_n t)
\end{equation*}
\begin{equation*}
  \ddot{x}(t) = -\omega_n^{2}\,x(t)
\end{equation*}

Constantes con condiciones iniciales
$x(0)=x_0$, $\dot{x}(0)=v_0$:
\begin{equation*}
  A = x_0,
  \qquad
  B = \frac{v_0}{\omega_n}
\end{equation*}

Forma equivalente amplitud--fase:
\begin{equation*}
  x(t) = X\sin(\omega_n t + \phi)
\end{equation*}
\begin{equation*}
  \dot{x}(t) = X\,\omega_n\cos(\omega_n t + \phi)
\end{equation*}
\begin{equation*}
  \ddot{x}(t) = -X\,\omega_n^{2}\sin(\omega_n t + \phi)
\end{equation*}
\begin{equation*}
  X = \sqrt{A^{2}+B^{2}} = \sqrt{x_0^{2}+\left(\frac{v_0}{\omega_n}\right)^{2}}
\end{equation*}
\begin{equation*}
  \phi = \arctan\!\left(\frac{A}{B}\right)
       = \arctan\!\left(\frac{x_0\,\omega_n}{v_0}\right)
\end{equation*}

En sistemas rotacionales se sustituye $x \rightarrow \theta$,
$m \rightarrow I_0$, $k \rightarrow k_t$.

Periodo y frecuencia:
\begin{equation*}
  T = \frac{2\pi}{\omega_n}\ [\mathrm{s}],
  \qquad
  f = \frac{1}{T} = \frac{\omega_n}{2\pi}\ [\mathrm{Hz}]
\end{equation*}

\seccion{Frecuencias naturales}

Sistema masa--resorte:
\begin{equation*}
  \omega_n = \sqrt{\frac{k}{m}}
\end{equation*}

Péndulo simple:
\begin{equation*}
  \omega_n = \sqrt{\frac{g}{l}}
\end{equation*}

Péndulo físico (compuesto):
\begin{equation*}
  \omega_n = \sqrt{\frac{m\,g\,l}{I_0}}
\end{equation*}

Péndulo de torsión:
\begin{equation*}
  \omega_n = \sqrt{\frac{k_t}{I_0}}
\end{equation*}

Teorema de ejes paralelos (Steiner):
\begin{equation*}
  I_0 = I_{cm} + m\,l^{2}
\end{equation*}

\seccion{Rigidez equivalente}

Resortes en serie:
\begin{equation*}
  \frac{1}{k_{eq}} = \sum_{i=1}^{n}\frac{1}{k_i}
  \qquad
  k_{eq} = \frac{1}{\displaystyle\sum_{i=1}^{n}\frac{1}{k_i}}
\end{equation*}

Resortes en paralelo:
\begin{equation*}
  k_{eq} = \sum_{i=1}^{n} k_i
\end{equation*}

Barra a tracción/compresión:
\begin{equation*}
  k = \frac{E\,A}{l}
\end{equation*}

Eje a torsión:
\begin{equation*}
  k_t = \frac{J_0\,G}{l}
\end{equation*}

Cuerpo sobre plano inclinado:
\begin{equation*}
  k_{eq} = \frac{M\,g\,\sin\alpha}{2\,R\,\theta}
\end{equation*}

Viga en voladizo con carga en el extremo:
\begin{equation*}
  y_{max} = \frac{P\,L^{3}}{3\,E\,I},
  \qquad
  k = \frac{P}{y_{max}} = \frac{3\,E\,I}{L^{3}}
\end{equation*}

\seccion{Propiedades del material}

Módulo de elasticidad (Young):
\begin{equation*}
  E = \frac{\sigma}{\varepsilon} = \frac{P\,L}{A\,\Delta L}
\end{equation*}

Esfuerzo normal y deformación unitaria:
\begin{equation*}
  \sigma = \frac{P}{A},
  \qquad
  \varepsilon = \frac{\Delta L}{L}
\end{equation*}

Esfuerzo cortante por torsión:
\begin{equation*}
  \tau = \frac{M\,R}{J_0}
\end{equation*}

Deformación angular por cortante:
\begin{equation*}
  \gamma = \frac{R\,\theta}{L}
\end{equation*}

Módulo de rigidez al cortante:
\begin{equation*}
  G = \frac{\tau}{\gamma}
\end{equation*}

Momento polar de inercia (sección circular maciza):
\begin{equation*}
  J_0 = \frac{\pi d^{4}}{32} = \frac{\pi R^{4}}{2}
\end{equation*}

Momento de inercia de área:
\begin{equation*}
  I_{rect} = \frac{b\,h^{3}}{12},
  \qquad
  I_{circ} = \frac{\pi d^{4}}{64}
\end{equation*}

\seccion{Nomenclatura}

\begin{tabbing}
\hspace{1.3cm}\=\kill
$x$ \> posición [m] \\
$\dot{x}$ \> velocidad [m/s] \\
$\ddot{x}$ \> aceleración [m/s$^2$] \\
$A,B$ \> constantes de integración [m] \\
$X$ \> amplitud [m] \\
$\phi$ \> ángulo de fase [rad] \\
$\omega_n$ \> frecuencia natural [rad/s] \\
$T$ \> periodo [s] \\
$f$ \> frecuencia [Hz] \\
$m,\,M$ \> masa [kg] \\
$k$ \> rigidez [N/m] \\
$k_t$ \> rigidez torsional [N$\cdot$m/rad] \\
$E$ \> módulo de Young [Pa] \\
$G$ \> módulo cortante [Pa] \\
$A$ \> área de sección [m$^2$] \\
$l,\,L$ \> longitud [m] \\
$I_0$ \> inercia de masa [kg$\cdot$m$^2$] \\
$I$ \> inercia de área [m$^4$] \\
$J_0$ \> inercia polar de área [m$^4$] \\
$P$ \> carga puntual [N] \\
$y_{max}$ \> deflexión máxima [m] \\
$b,h$ \> base y altura de sección [m] \\
$d$ \> diámetro [m] \\
$\sigma$ \> esfuerzo normal [Pa] \\
$\tau$ \> esfuerzo cortante [Pa] \\
$\varepsilon$ \> deformación unitaria [--] \\
$\gamma$ \> deformación cortante [--] \\
$\theta$ \> ángulo de giro [rad] \\
$g$ \> $9{,}81$ m/s$^2$ \\
\end{tabbing}

\end{multicols}
\end{document}
